
The generation of Ohmic current is based on the concept of transformer,
where the plasma torus is the secondary circuit. An electric field is
induced by the temporal variation $\partial \psi /\partial t$ of the
poloidal flux generated by the primary circuit. The induced current density
is then calculated by Ohm's law, $\mathbf{J}=\Bbb{\sigma }_{\Omega }\mathbf{E%
},$ where $\Bbb{\sigma }_{\Omega }$ is the electric conductivity calculated
by accounting for the Coulomb interaction between the strongly magnetized
components of the plasma. Using Faraday's law, 
\[
\frac{\partial \mathbf{B}}{\partial t}=-\nabla \times \mathbf{E} 
\]

We consider the surface $S\left( \psi ,\theta \right) $ which is a truncated
cone delimited by two circles being $C_{1}$, the magnetic axis, and $%
C_{2}\left( \psi ,\theta \right) $, the toroidal line at position $\left(
\psi ,\theta \right) $. Applying the integral formula to Faraday's law, we
have 
\begin{equation}
\iint_{S\left( \psi ,\theta \right) }\mathbf{dS}\cdot \frac{\partial \mathbf{%
B}}{\partial t}=-\iint_{S\left( \psi ,\theta \right) }\mathbf{dS}\cdot
\nabla \times \mathbf{E=}-\oint_{C\left( \psi ,\theta \right) }\mathbf{dl}%
\cdot \mathbf{E}
\end{equation}

The poloidal flux is given by 
\begin{equation}
\psi =\frac{1}{2\pi }\iint_{S\left( \psi ,\theta \right) }dSB_{P}=\frac{1}{2%
\pi }\iint_{S\left( \psi ,\theta \right) }\mathbf{dS\cdot B}
\end{equation}
so that we get 
\begin{eqnarray}
2\pi \frac{\partial \psi }{\partial t} &=&\oint_{C_{1}}\mathbf{dl}\cdot 
\mathbf{E}-\oint_{C_{2}\left( \psi ,\theta \right) }\mathbf{dl}\cdot \mathbf{%
E} \\
&=&\int_{0}^{2\pi }R_{p}d\phi E_{\phi 0}-\int_{0}^{2\pi }R\left( \psi
,\theta \right) d\phi E_{\phi }\left( \psi ,\theta \right) 
\end{eqnarray}
where $R_{p}$ is the major radius on axis and $E_{\phi p}$ is the electric
field on axis. Using axisymmetry, this becomes 
\begin{equation}
\frac{\partial \psi }{\partial t}=R_{p}E_{\phi p}-R\left( \psi ,\theta
\right) E_{\phi }\left( \psi ,\theta \right) 
\end{equation}
and therefore 
\begin{equation}
RE_{\phi }=R_{p}E_{\phi p}-\frac{\partial \psi }{\partial t}
\end{equation}
and $RE_{\phi }$ is only a function of $\psi $. We can therefore rewrite 
\begin{equation}
E_{\phi }\left( \psi ,\theta \right) =E_{\phi 0}\frac{R_{0}}{R}
\end{equation}
where $R_{0}$ is the major radius taken at the the poloidal position $\theta
_{0}\ $where the magnetic field $\mathbf{B}$ is minimum on a flux-surface.By
deifinition,
\begin{equation}
E_{\Vert 0}\left( \psi \right) =E_{\Vert }\left( \psi ,\theta _{0}\right) 
\end{equation}

The electric field along the field line can then be obtained by projection,
which gives 
\begin{equation}
E_{\Vert }=\left( \widehat{b}\cdot \widehat{\phi }\right) E_{\phi }=\frac{%
B_{T}}{B}E_{\phi }
\end{equation}
so that we get 
\begin{equation}
\frac{B}{B_{T}}E_{\Vert }=\frac{B_{0}}{B_{T0}}E_{\Vert 0}\frac{R_{0}}{R}
\end{equation}
and then 
\begin{equation}
E_{\Vert }=\frac{B_{T}}{B_{T0}}\frac{B_{0}}{B}\frac{R_{0}}{R}E_{\Vert
0}\left( \psi \right) =\frac{1}{\Psi }\frac{R_{0}^{2}}{R^{2}}E_{\Vert
0}\left( \psi \right) \label{efieldpar}
\end{equation}
with $\Psi =B/B_{0}$ as defined in Sec. \ref{NormBounceSection}.

\subsection{Conservative Form for the Ohmic Electric Field Operator}

The effect of the electric fied $E_{\Vert }$ can be expressed in a
conservative form as $\mathcal{E}(f_{0})=\nabla _{\mathbf{p}}\cdot \mathbf{S}%
^{\text{E}}$, where the flux in momentum space is easily expressed in
cylindrical coordinates $\left( p_{_{\Vert }},p_{\bot }\right) $ as 
\begin{equation}
\mathbf{S}_{\left( p_{\Vert },p_{\bot }\right) }^{\text{E}}=\left( 
\begin{array}{l}
S_{p_{_{\Vert }}}^{E} \\ 
S_{p_{\bot }}^{E}
\end{array}
\right)
\end{equation}
with 
\begin{eqnarray}
S_{p_{\Vert }}^{\text{E}} &=&q_{e}E_{\Vert }f_{0} \\
S_{p_{\bot }}^{\text{E}} &=&0
\end{eqnarray}
where $q_{e}$ is the electronic charge.

The transformation from cylindrical to spherical coordinates is given by 
\begin{equation}
\mathbf{S}_{\left( p,\xi \right) }^{\text{E}}=\mathcal{R}^{-1}\mathbf{S}%
_{\left( p_{\Vert },p_{\bot }\right) }^{\text{E}}
\end{equation}
where $\mathcal{R}$ is the rotational matrix 
\[
\mathcal{R}=\left( 
\begin{array}{ll}
\xi & -\sqrt{1-\xi ^{2}} \\ 
\sqrt{1-\xi ^{2}} & \xi
\end{array}
\right) 
\]

Using $\mathcal{R}^{-1}=$ $^{t}\mathcal{R}$ we find 
\begin{eqnarray}
S_{p}^{\text{E}} &=&\xi S_{p_{\Vert }}^{\text{E}} \\
S_{\xi }^{\text{E}} &=&-\sqrt{1-\xi ^{2}}S_{p_{\Vert }}^{\text{E}}
\end{eqnarray}
and $\mathbf{S}^{\text{E}}$ countains a convective part only 
\begin{equation}
\mathbf{S}^{\text{E}}=\mathbf{F}_{\mathbf{p}}^{\text{E}}f_{0}  \label{elec}
\end{equation}
with 
\begin{eqnarray*}
F_{p}^{\text{E}} &=&q_{e}\xi E_{\Vert } \\
F_{\xi }^{\text{E}} &=&-q_{e}\sqrt{1-\xi ^{2}}E_{\Vert }
\end{eqnarray*}
